\section{Discussions}
\label{asec:discussions}


\subsection{Limitations}
\label{assec:limitations}

While comprehensive with 26K queries, \dataset represents only a snapshot of the vast space of possible open-ended queries and may not capture all forms of creative divergence across different contexts. Moreover, the focus on English-language prompts derived from WildChat potentially underrepresents linguistic, cultural, and regional diversity in user interactions with language models, limiting the generalizability of our findings to non-English contexts and diverse cultural perspectives. To construct our taxonomy of queries, we used GPT-4o to efficiently label user queries, and we found human annotators achieved 74.7\% agreement with these automatic labels. While this accuracy could be improved, e.g., by labeling with stronger models or through model ensembling approaches, this performance is comparable to typical human annotation accuracy and we deem sufficient to represent the distribution of categories within our open-ended taxonomy.
Finally, while our analysis successfully reveals clear patterns such as the ``time is a river/weaver'' dichotomy, it may oversimplify the multidimensional nature of creative expression, and relying on semantic similarity of text embeddings to quantify diversity may lack sufficient expressiveness to capture the full spectrum of creative variation in generated responses.

While extending \dataset to multilingual and multicultural settings is an important direction, we anticipate that similar homogenization issues will likely arise across languages and cultures, given the global overlap in pretraining data sources and alignment practices. We encourage future work to pursue such extensions, and we hope our benchmark provides a comprehensive foundation that multilingual research can build upon. Our taxonomy was designed to be language-agnostic and can support adaptation when appropriate resources and expertise are available. We'll further strengthen this discussion in camera-ready.

Diversity and quality represent two key dimensions in evaluating responses to open-ended queries. In this work, we focused primarily on diversity, with fixed model decoding configurations, without studying quality. We selected a decoding configuration that generated empirically coherent text and investigated whether models can be guided to produce responses that maximize diversity while maintaining reasonable quality. Furthermore, we show that existing LMs, LM judges, and reward models struggle to reliably distinguish between equal-quality responses for open-ended queries. This motivates the need for further investigation into reliable automatic evaluation of equal-quality responses for open-ended queries, beyond human annotation. Finally, although our analysis identifies clear evidence of an ``Artificial Hivemind'' effect across models, it falls short of establishing the underlying mechanisms causing this homogenization. Future research can mechanistically disentangle causality, whether stemming from shared training data, memorization, generalization, or other factors.

Our work has significant implications for AI development and society, providing a valuable dataset for diagnosing and potentially mitigating LM homogenization and accelerating the development of more diverse, creative AI systems that align with diverse human needs and perspectives. However, this ``Artificial Hivemind'' effect also raises serious concerns about models' long-term impact on human creativity. If users increasingly rely on such systems for creative tasks, exposure to homogenized outputs could subtly influence human thinking patterns and reduce overall cultural and intellectual diversity. Furthermore, it could exacerbate existing biases and limit representation of marginalized perspectives. When LMs converge on dominant cultural expressions---such as Western-centric metaphors like ``time is a river"---they may inadvertently suppress alternative worldviews and traditions.

The ``Artificial Hivemind'' effect also highlights fundamental questions about what values we want AI systems to embody. Should AI prioritize efficiency and consistency, or diversity and novelty? These choices reflect deeper societal values about creativity, culture, and human flourishing. The ethical implications of LM homogenization extend far beyond technical considerations, touching on fundamental questions about human agency, cultural preservation, and the kind of future we want to create with these technologies. Addressing these challenges requires not just technical innovation, but sustained ethical reflection and inclusive dialogue about the role of AI in human creative expression.


\subsection{Future Directions}
\label{assec:future}

Our findings open several promising directions for advancing the study and mitigation of behavioral convergence in large language models.

\textbf{Foundation and training analysis.} We plan to extend the Artificial Hivemind testbed to foundation models without instruction-following capabilities to better disentangle the respective roles of pre-training and post-training in shaping convergent behaviors. Moreover, we aim to quantify the relative contributions of different post-training pipelines—such as supervised fine-tuning, RLHF/RLAIF, and constitutional training—to the emergence of homogenized responses.

\textbf{Mitigation and alignment strategies.} Future work will explore diversity-aware training objectives and alignment schemes that explicitly reward exploration of multiple valid modes while preserving response quality. We also intend to benchmark decoding strategies (e.g., diverse beam search, nucleus sampling variants) under the Artificial Hivemind metric to evaluate their effectiveness in counteracting homogenization.

\textbf{Practical integration.} To make our framework actionable, we will integrate Artificial Hivemind into red-teaming workflows to stress-test model robustness and coverage under open-ended queries. Additionally, our dataset can serve as a training prompt resource for reinforcement learning methods that explicitly encourage diversity during alignment. Finally, Hivemind’s diagnostic signals can inform curriculum design—gradually exposing models to increasingly open-ended prompts that are most susceptible to mode collapse.

Together, these directions position Artificial Hivemind as not only a descriptive diagnostic tool, but also a foundation for developing training, decoding, and evaluation practices that foster greater diversity and individuality in language models.

\subsection{Broader Implications}
\label{assec:broad}

\paragraph{Societal Implications.}
While some degree of convergence among large language models (LLMs) is statistically expected, its societal consequences warrant close scrutiny. As billions of users increasingly depend on LLMs for creative, educational, and decision-making purposes, understanding and quantifying behavioral homogenization becomes critical. Emerging evidence shows measurable shifts in human writing styles, creative ideation, and divergent thinking following the widespread adoption of systems like ChatGPT \cite{Kumar_2025,Anderson2024HomogenizationEO,Ashkinaze_2025,doi:10.1126/sciadv.adn5290,kosmyna2025brainchatgptaccumulationcognitive,sourati2025shrinkinglandscapelinguisticdiversity}. These findings suggest that model-level convergence may propagate into human expression, amplifying uniformity in linguistic and cognitive patterns at scale. By introducing quantitative tools to measure and track such convergence, our work provides an empirical foundation for assessing the long-term cultural and epistemic consequences of LLM-mediated communication.

\paragraph{Implications for Data Distillation and Model Training.}
Our results also have immediate implications for synthetic data generation and model training practices. While distilling knowledge from LLMs has proven transformative for efficiency and scalability, it is well known that relying on a single model as a teacher can intensify mode collapse, reinforcing narrow response patterns, diminishing output diversity, and, in extreme cases, leading to degenerative feedback loops \cite{alemohammad2024selfconsuming,kazdan2025collapsethriveperilspromises}. To counter this, recent approaches employ model swarms and multi-agent frameworks that aggregate outputs from multiple models in pursuit of diversity \cite{feng2025modelswarmscollaborativesearch,jiang2025spartaalignmentcollectivelyaligning}.


However, our findings reveal a fundamental limitation: even across distinct state-of-the-art models, diversity in open-ended tasks is far from guaranteed. Models often converge toward highly similar answers, undermining the assumed benefits of multi-model distillation. This insight carries significant implications for research on synthetic data curation, ensemble-based alignment, and constitutional AI. Without rigorous diagnostic frameworks such as Artificial Hivemind, practitioners risk overestimating cross-model diversity and unintentionally perpetuating homogenization across generations of language models.

